\subsection{From the source kernel bounds to an actual free coefficient basis}
\label{subsec:faithful-kernel-norm-free}

This is the algebraic producer
\srcid{KernelRankAlgebraicBound}; its conclusion is a basis, not
merely a dimension assertion.

\begin{srproposition}\label{prop:faithful-kernel-norm-free}
Let \(E\) be a number field, \(O=\mathcal O_E\), and
\(b_1,\ldots,b_n\) its specified integral basis, where
\(n=[E:\mathbb Q]\).  Let \(k\) be algebraically closed of
characteristic the prime \(p\), let \(W\) be a finite-dimensional
\(k\)-vector space, and let
\(\rho:O\to\operatorname{End}_k(W)\) be a unital ring homomorphism.
Fix \(r\ge0\).  Assume, for every \(a\in O\), that
\[
 \operatorname{charpol}(\rho(a))
  =\left(\operatorname{charpol}(m_a:O\to O)
                                \bmod k\right)^r,
 \tag{NF1}\label{eq:faithful-norm-free-charpoly}
\]
and, for every \(a\ne0\), that
\[
       \dim_k\ker\rho(a)\le
                       r\,v_p(|N_{E/\mathbb Q}(a)|).
 \tag{NF2}\label{eq:faithful-norm-free-kernel}
\]
Then there exist \(x_1,\ldots,x_r\in W\) such that
\[
        \{\rho(b_j)x_l:1\le j\le n,\ 1\le l\le r\}
 \tag{NF3}\label{eq:faithful-norm-free-basis}
\]
is a \(k\)-basis of \(W\).  Equivalently the extended action makes
\(W\) free of rank \(r\) over \(k\otimes_{\mathbb Z}O\).
\end{srproposition}

\begin{proof}
Put \(A=k\otimes_{\mathbb Z}O\).  The rule
\((c\otimes a)w=c\rho(a)w\) gives the actual \(A\)-module
structure.  For \(a\in O\), the regular multiplication
characteristic polynomial on \(A\) is the image in \(k[T]\) of
the integral multiplication characteristic polynomial.
Thus \eqref{eq:faithful-norm-free-charpoly} compares the
characteristic polynomials on \(W\) and on the regular module
\(A\) for each element \(1\otimes a\).
Taking degrees gives
\[
                         \dim_kW=nr.
 \tag{NF4}\label{eq:faithful-norm-free-dimension}
\]
We keep the source's local-generator argument.

\emph{The two arithmetic elements.}
For a maximal ideal \(\mathfrak m\subset A\), its residue field
is \(k\), since it is finite-dimensional over the algebraically
closed field \(k\).  Choose the corresponding \(k\)-algebra map
\(\chi_0:A\to k\).
The restriction \(\psi:O\to k\), \(a\mapsto\chi_0(1\otimes a)\),
has a nonzero maximal kernel \(\mathfrak w\) containing \(p\).
Write
\[
                      |O/\mathfrak w|=p^f .
\]
The source constructs a nonzero \(\pi\in\mathfrak w\) with
\[
                       v_p(|N\pi|)=f .
 \tag{NF5}\label{eq:faithful-norm-free-pi}
\]
Here is its ideal calculation.
The Dedekind two-generator statement gives \(\pi\) such that
\(\mathfrak w(p)+(\pi)=\mathfrak w\).
Consequently \((\pi)=\mathfrak w I\) and cancellation of the
nonzero ideal \(\mathfrak w\) gives \((p)+I=O\).
The ideal \(I\) is nonzero: otherwise \(p\) would be a unit
in \(O\), contrary to \(|N(p)|=p^n>1\).
Thus \(O/I\) is finite, and multiplication by \(p\) in that
quotient is invertible.  Hence \(p\nmid |O/I|\): otherwise Cauchy's theorem
would give a nonzero element of additive order \(p\), contradicted
by that invertibility.
Taking ideal norms gives
\(|N\pi|=p^f|O/I|\), proving
\eqref{eq:faithful-norm-free-pi}, including \(\pi\ne0\).

Factor
\[
                  (p)=\mathfrak w^c Q,\qquad
                  \mathfrak w+Q=O.
\]
Choose \(z\in Q\) with \(z\equiv1\bmod\mathfrak w\), choose a
lift \(g_0\in O\) of a generator of the finite cyclic group
\((O/\mathfrak w)^\times\), and put \(\theta=g_0z\).
Its residue is nonzero and generates that multiplicative group.
For every \(\pi'\in\mathfrak w\),
\[
                         (\pi'\theta)^{c+1}\in(p).
 \tag{NF6}\label{eq:faithful-norm-free-theta}
\]
Indeed one factor contributes membership in \(Q\), while \(c\)
factors of \(\pi'\) contribute membership in \(\mathfrak w^c\).

There are exactly \(f\) embeddings
\(j:O/\mathfrak w\to k\); finite-field separability gives this
count.  Extend them to characters
\[
 \chi_j:A\to k,\qquad
                  \chi_j(c\otimes a)=c\,j(a\bmod\mathfrak w).
\]
For
\[
                 t=1\otimes\theta,\qquad u=1\otimes\pi
\]
the \(f\) numbers \(\alpha_j=\chi_j(t)\) are pairwise distinct
and nonzero: equality on the chosen multiplicative generator
gives equality of the embeddings.
Equation \eqref{eq:faithful-norm-free-theta} says \(ut\) is
nilpotent in \(A\).  Moreover
\[
 \dim_k\ker(u:W\to W)\le rf
 \tag{NF7}\label{eq:faithful-norm-free-total-kernel}
\]
by \eqref{eq:faithful-norm-free-kernel} and
\eqref{eq:faithful-norm-free-pi}.

\emph{The generalized eigenspaces.}
For each \(j\), let \(G_j(A)\) and \(G_j(W)\) be the maximal
generalized \(\alpha_j\)-eigenspaces for multiplication by \(t\).
Set \(d_j=\dim_kG_j(A)\).
The characteristic-polynomial assumption at \(\theta\) gives
\[
                      \dim_kG_j(W)=r d_j .
 \tag{NF8}\label{eq:faithful-norm-free-eigenspace-size}
\]
We justify both \(d_j>0\) and the nilpotence exponent used next.

For any \(v\) in an Artinian commutative ring, stabilization of
its principal powers supplies \(n,y\) with
\(v^{n+1}y=v^n\).
Then
\[
             e=1-(vy)^n
\]
satisfies \(v^ne=0\), and \(1-e\) is divisible by every power
\(v^M\).  For the latter assertion use
\(v^n(vy)^M=v^n\); explicitly
\(1-e=v^{n+M}y^{n+M}\).
This is the source's \srcid{fitting} construction.
Apply it to \(v=t-\alpha_j\).
The resulting \(e_j\) lies in \(G_j(A)\), and
\(\chi_j(e_j)=1\), since \(1-e_j\) is divisible by \(t-\alpha_j\).
Thus \(d_j>0\).
It also acts as the identity on \(G_j(W)\):
for any such vector \(w\), choose \(M\) with
\((t-\alpha_j)^Mw=0\) and use the divisibility of \(1-e_j\).

Since \(\alpha_j\ne0\), the elements \(t\) and \(t-\alpha_j\)
are coprime.  Because \(ut\) is nilpotent, Bezout applied to
powers of those two elements shows that \(u\) is nilpotent
on \(G_j(A)\).  More explicitly, if \((ut)^N=0\) and
\((t-\alpha_j)^Mv=0\), write
\(a t^N+b(t-\alpha_j)^M=1\) and multiply by \(u^Nv\).
Both terms on the right vanish.
Cayley--Hamilton on this \(d_j\)-dimensional subspace gives
\[
                          u^{d_j}G_j(A)=0 .
\]
In particular \(u^{d_j}e_j=0\).
As \(e_j\) acts as the identity on \(G_j(W)\), we obtain
\[
                          u^{d_j}G_j(W)=0 .
 \tag{NF9}\label{eq:faithful-norm-free-nilpotence}
\]

For a nilpotent endomorphism \(U\) with \(U^d=0\), the elementary
source bound is
\[
                     \dim V\le d\,\dim\ker U.
 \tag{NF10}\label{eq:faithful-nilpotent-corank}
\]
To see its exact mechanism, lift a basis of \(V/UV\), whose
dimension is \(\dim\ker U\), to \(y_1,\ldots,y_s\).
Then \(V=\operatorname{span}_k(y_i)+UV\).
Iteration shows that \(U^q y_i\), \(0\le q<d\), span \(V\).
This proves \eqref{eq:faithful-nilpotent-corank} without assuming
a Jordan basis.
Apply it to \(u|G_j(W)\), using
\eqref{eq:faithful-norm-free-eigenspace-size} and
\eqref{eq:faithful-norm-free-nilpotence}.
As \(d_j>0\),
\[
                          \dim_k\ker(u|G_j(W))\ge r.
\]
These kernel spaces lie in pairwise independent generalized
eigenspaces, all within \(\ker(u|W)\).
There are \(f\) of them.  Equation
\eqref{eq:faithful-norm-free-total-kernel} therefore forces
\[
                       \dim_k\ker(u|G_j(W))=r
                           \quad\text{for every }j.
 \tag{NF11}\label{eq:faithful-norm-free-local-corank}
\]

\emph{Generators modulo the chosen maximal ideal.}
For the embedding inducing \(\chi_0\), use
\eqref{eq:faithful-norm-free-local-corank} to lift a basis of
\(G_j(W)/uG_j(W)\) to \(x_1,\ldots,x_r\in G_j(W)\).
It follows that
\[
       G_j(W)=\operatorname{span}_k(x_1,\ldots,x_r)+uG_j(W).
\]
For arbitrary \(w\in W\), the vector \(e_jw\) is in \(G_j(W)\).
The two remaining terms in the decomposition
\[
          w=e_jw+(1-e_j)w
\]
are controlled as follows.
We have \(u\in\ker\chi_0=\mathfrak m\), because \(ut\) is
nilpotent and \(\chi_0(t)\ne0\); also \(1-e_j\in\mathfrak m\)
by its divisibility by \(t-\alpha_j\).
Hence
\[
 W=A x_1+\cdots+A x_r+\mathfrak mW .
 \tag{NF12}\label{eq:faithful-norm-free-local-generators}
\]
This is the actual \srcid{core}/\srcid{locmax} conclusion.

\emph{Gluing the generators.}
The Artinian algebra \(A\) has finitely many maximal ideals.
For each \(\mathfrak m\), choose \(r\) vectors
\(y_{\mathfrak m,l}\) as in
\eqref{eq:faithful-norm-free-local-generators}.
The Chinese remainder theorem gives \(c_{\mathfrak m}\in A\)
whose residue is \(1\) at \(\mathfrak m\) and \(0\) at all
other maximal ideals.
Set
\[
               x_l=\sum_{\mathfrak m}
                                  c_{\mathfrak m}y_{\mathfrak m,l},
       \qquad N=\sum_l A x_l.
\]
Then \(W=N+\mathfrak mW\) for every maximal \(\mathfrak m\).
The source's determinant trick proves \(N=W\).
Indeed, if the ideal
\(J=\{a\in A:aW\subset N\}\) were proper, choose
\(\mathfrak m\supset J\).
The finitely generated equality \(W=N+\mathfrak mW\) gives
an element \(a\equiv1\bmod\mathfrak m\) with \(aW\subset N\):
write finitely many module generators in terms of one another
with coefficients in \(\mathfrak m\), and apply the adjugate
to the matrix \(I-M\).
Thus \(a\in J\subset\mathfrak m\), a contradiction.
Consequently \(1\in J\) and \(W=N\).

The \(A\)-generators \(x_l\) imply that the vectors in
\eqref{eq:faithful-norm-free-basis} span over \(k\):
expand an arbitrary element of \(O\) in the integral basis
\(b_j\), then expand a tensor in \(A\) as a sum of pure tensors.
There are \(nr\) such vectors, exactly \(\dim_k W\) by
\eqref{eq:faithful-norm-free-dimension}; they are therefore
linearly independent.  This proves the desired actual basis.
\end{proof}

\paragraph{Frozen-source ledger.}
The complete source is
\href{https://github.com/ChristianKudeba/alpoge-rh-verification/blob/alpoge-source-7cbf7a0c/src/Development/WeilBound/CoreStatements/TorusUniformisation/ChowMahlerFamilies/KernelRankAlgebraicBound.lean}
{\srcid{KernelRankAlgebraicBound}}
(blob \srcid{ed4a911401a5dc7a2c756d59035b64c06fb41f13}).
The paragraphs above follow, in order, its
\srcid{arith\_pi}, \srcid{arith\_theta}, \srcid{count\_hom},
\srcid{fitting}, \srcid{pow\_finrank\_kill},
\srcid{finrank\_G}, \srcid{core}, \srcid{locmax},
\srcid{glue}, and final theorem.
The nilpotent spanning argument is its imported
\href{https://github.com/ChristianKudeba/alpoge-rh-verification/blob/alpoge-source-7cbf7a0c/src/Development/WeilBound/CoreStatements/TorusUniformisation/ChowMahlerFamilies/JordanBasisNilpotent.lean}
{\srcid{JordanBasisNilpotent}},
\srcid{exists\_span\_sup\_range} and \srcid{span\_eq\_top};
its construction is given explicitly in
\eqref{eq:faithful-nilpotent-corank}.
The wrapper
\href{https://github.com/ChristianKudeba/alpoge-rh-verification/blob/alpoge-source-7cbf7a0c/src/Development/AutomorphicSpectrum/Proofs/AlgebraProofs/KernelRankBoundNormValuation.lean}
{\srcid{KernelRankBoundNormValuation.lemF\_algebra}}
has exactly the hypotheses and conclusion of this proposition.

This is a closed algebraic implication.
To apply it to the geometric doubled action one must still
produce \eqref{eq:faithful-norm-free-kernel}, with the same
\(r\) as the actual characteristic-polynomial identity.
The proper-germ divisibility proves one component of that
kernel bound.  The torus lattice, quotient gauge, and
norm-degree calculation supply the other components; the
present proposition does not replace any of them.
